\chapter{Hermitian Forms}\label{ch:b2:hermitian}

Complex vector spaces have their own inner-product geometry, with
one twist: linearity in one variable, \emph{conjugate}-linearity in
the other. The payoff for accepting the twist is a spectral theory
even cleaner than the real one --- Hermitian endomorphisms have real
eigenvalues, unitary ones have unimodular eigenvalues, and both
diagonalize in orthonormal bases. This short chapter runs the
Euclidean program of \cref{ch:b2:quadratic} over $\C$.

\section{Hermitian inner products}

\begin{definition}\label{def:b2:hermitian:def}
A \emph{Hermitian inner product}\index{Hermitian inner product} on
a complex vector space $E$ is a map $\langle\cdot,\cdot\rangle
\colon E \times E \to \C$ that is linear in the second variable,
conjugate-symmetric ($\langle y, x\rangle = \conj{\langle x,
y\rangle}$ --- hence conjugate-linear in the first variable), and
positive definite ($\langle x, x\rangle > 0$ for $x \neq 0$). The
standard example on $\C^n$:
\[
\langle x, y \rangle = \sum_{i=1}^{n} \conj{x_i}\, y_i ;
\]
on continuous functions, $\langle f, g\rangle = \int_a^b \conj
f\,g$. Norm: $\norm x = \sqrt{\langle x,x\rangle}$; a
finite-dimensional complex space so equipped is a \emph{Hermitian
space}.
\end{definition}

\begin{example}[First computations]\label{ex:b2:hermitian:firstcomp}
In $\C^2$, take $x = (1+\iu,\ 2-\iu)$ and $y = (\iu,\ 1)$.
Then
\[
\norm x^2 = \abs{1+\iu}^2 + \abs{2-\iu}^2 = 2 + 5 = 7,
\qquad
\norm y^2 = 1 + 1 = 2,
\]
and, conjugating the first argument,
\[
\langle x, y\rangle
= \conj{(1+\iu)}\,\iu + \conj{(2-\iu)}\cdot1
= (1-\iu)\iu + (2+\iu)
= (\iu + 1) + (2 + \iu) = 3 + 2\iu .
\]
Cauchy--Schwarz checks out: $\abs{\langle x, y\rangle}^2 = 9 +
4 = 13 \leq 14 = \norm x^2\norm y^2$ --- close to equality,
because $x$ is close to a multiple of $y$. Note also
$\langle y, x\rangle = \conj{3 + 2\iu} = 3 - 2\iu$:
conjugate-symmetry in action, and the reason $\langle x,
x\rangle$ is always real.
\end{example}

\begin{theorem}[Cauchy--Schwarz, complex]\label{thm:b2:hermitian:cs}
$\abs{\langle x, y\rangle} \leq \norm x \norm y$, with equality iff
$x, y$ are linearly dependent; $\norm\cdot$ is a norm. Moreover
orthonormal bases exist (Gram--Schmidt runs verbatim), with
\[
x = \sum_i \langle e_i, x\rangle\, e_i,
\qquad
\norm x^2 = \sum_i \abs{\langle e_i, x\rangle}^2 .
\]
\end{theorem}

\begin{proof}
For $y \neq 0$ and $t \in \C$: $0 \leq \norm{x - ty}^2 = \norm x^2
- 2\Re\bigl(\conj t\langle y, x\rangle\bigr) + \abs t^2\norm y^2$.
Choose $t = \frac{\langle y, x\rangle}{\norm y^2}$:
\[
0 \leq \norm x^2 - \frac{\abs{\langle y, x\rangle}^2}{\norm y^2},
\]
which is the inequality; equality forces $x = ty$. Triangle
inequality, in full:
\[
\norm{x + y}^2 = \norm x^2 + 2\,\Re\langle x, y\rangle
+ \norm y^2
\leq \norm x^2 + 2\,\abs{\langle x, y\rangle} + \norm y^2
\leq \bigl(\norm x + \norm y\bigr)^2 ,
\]
using $\Re z \leq \abs z$ and then Cauchy--Schwarz; homogeneity
and separation are immediate, so $\norm\cdot$ is a norm.
Gram--Schmidt: as in the real case, with conjugates
placed by the definition (note the convention: our products are
conjugate-linear in the \emph{first} slot, so coordinates are
$\langle e_i, x\rangle$; the next example runs the algorithm
once in full).
\end{proof}

\begin{example}[Complex Gram--Schmidt, run in full]\label{ex:b2:hermitian:gramschmidt}
Orthonormalize the basis $v_1 = (1, \iu)$, $v_2 = (0, 1)$ of
$\C^2$. First vector: $\norm{v_1}^2 = \abs1^2 + \abs\iu^2 = 2$,
so $e_1 = \frac{1}{\sqrt2}(1, \iu)$. Project $v_2$ --- with the
conjugate in the first slot:
\[
\langle e_1, v_2\rangle
= \frac{1}{\sqrt2}\bigl(\conj{1}\cdot0 +
\conj{\iu}\cdot1\bigr)
= \frac{-\iu}{\sqrt2} ,
\qquad
v_2 - \langle e_1, v_2\rangle e_1
= (0,1) + \frac{\iu}{2}\,(1, \iu)
= \Bigl(\frac\iu2,\ \frac12\Bigr) .
\]
Its norm is $\sqrt{\frac14 + \frac14} = \frac{1}{\sqrt2}$:
$e_2 = \frac{1}{\sqrt2}(\iu, 1)$. Check: $\langle e_1,
e_2\rangle = \frac12(\conj1\cdot\iu + \conj\iu\cdot1) =
\frac12(\iu - \iu) = 0$. Coordinates of $v_2$ in the new basis:
$v_2 = \langle e_1, v_2\rangle e_1 + \langle e_2, v_2\rangle
e_2$ with $\langle e_2, v_2\rangle = \frac{1}{\sqrt2}$ ---
mind the order: $\langle v_2, e_1\rangle$ would give the
\emph{conjugate} coefficient. Closing insight: the algorithm is
the Euclidean one verbatim; the only trap is where the
conjugation falls, and computing $\norm{v_2 -
\text{proj}}^2 > 0$ silently uses positivity --- the axiom
that makes the whole geometry work.
\end{example}

\section{Adjoint, Hermitian and unitary endomorphisms}

\begin{definition}\label{def:b2:hermitian:adjoint}
The \emph{adjoint} $u^*$ of $u \in \mathcal{L}(E)$ is defined by
$\langle u^*(x), y\rangle = \langle x, u(y)\rangle$; in an
orthonormal basis, $\operatorname{Mat}(u^*) = \conj{A}^{\mathsf T} =:
A^{\dagger}$ (conjugate transpose) --- indeed, if $B =
(b_{ij})$ is the matrix of $u^*$ in the orthonormal basis
$(e_i)$, then $b_{ij} = \langle e_i, u^*(e_j)\rangle$, and the
defining identity gives
\[
\conj{b_{ij}}
= \langle u^*(e_j), e_i\rangle
= \langle e_j, u(e_i)\rangle = a_{ji} ,
\qquad\text{so}\qquad
B = \conj{A}^{\mathsf T} .
\] $u$ is
\emph{Hermitian}\index{Hermitian endomorphism} when $u^* = u$
($A^\dagger = A$), \emph{unitary}\index{unitary group} when $u^*u =
\mathrm{id}$ ($A^\dagger A = I$: the group $U(n)$), \emph{normal}
when $u^*u = uu^*$.
\end{definition}

\begin{proposition}\label{prop:b2:hermitian:eigenvalues}
Eigenvalues of a Hermitian endomorphism are \emph{real};
eigenvalues of a unitary endomorphism have \emph{modulus $1$}; in
both cases eigenspaces for distinct eigenvalues are orthogonal.
\end{proposition}

\begin{proof}
Hermitian, $u(x) = \lambda x$, $x \neq 0$:
\[
\lambda \norm x^2 = \langle x, u(x)\rangle = \langle u(x), x\rangle
= \conj{\langle x, u(x)\rangle} = \conj\lambda\,\norm x^2 ,
\]
so $\lambda \in \R$. Unitary: $\norm{u(x)} = \norm x$ (from $u^*u =
\mathrm{id}$), so $\abs\lambda\norm x = \norm x$. Orthogonality
(Hermitian case): $\lambda\langle x, y\rangle = \langle u(x),
y\rangle = \langle x, u(y)\rangle = \mu\langle x, y\rangle$ for
eigenvectors with real $\lambda \neq \mu$. Unitary case, in
full: for $u(x) = \lambda x$, $u(y) = \mu y$ with $\lambda \neq
\mu$ (both unimodular),
\[
\langle x, y\rangle = \langle u(x), u(y)\rangle
= \conj\lambda\mu\,\langle x, y\rangle ,
\]
and $\conj\lambda\mu = \frac{\mu}{\lambda} \neq 1$: the factor
is not $1$, so $\langle x, y\rangle = 0$.
\end{proof}

\begin{example}[A skew-Hermitian matrix, diagonalized]\label{ex:b2:hermitian:skewexample}
$A = \begin{pmatrix} 0 & -2\\ 2 & 0\end{pmatrix}$ satisfies
$A^\dagger = A^{\mathsf T} = -A$: skew-Hermitian (also real
antisymmetric --- over $\R$ it has no eigenvalues at all).
Characteristic polynomial $X^2 + 4$: eigenvalues $\pm2\iu$,
purely imaginary, as \cref{exo:b2:hermitian:9} predicts in
general. Eigenvectors: $(A - 2\iu I)v = 0$ gives $v_1 =
\frac{1}{\sqrt2}(1, \iu)$, and $v_2 = \frac{1}{\sqrt2}(1,
-\iu)$ for $-2\iu$; they are orthogonal:
\[
\langle v_1, v_2\rangle
= \tfrac12\bigl(\conj{1}\cdot1 +
\conj{\iu}\cdot(-\iu)\bigr)
= \tfrac12(1 - 1) = 0 .
\]
So $A = U\operatorname{diag}(2\iu, -2\iu)\,U^\dagger$ with $U =
(v_1\ v_2)$ unitary. Closing insight: $H = -\iu A =
\begin{pmatrix} 0 & 2\iu\\ -2\iu & 0\end{pmatrix}$ is Hermitian
with the real spectrum $\{\pm2\}$ and the \emph{same}
eigenvectors --- the bijection $u \mapsto \iu u$ between
Hermitian and skew-Hermitian endomorphisms
(\cref{exo:b2:hermitian:9}), seen matrix by matrix; over $\R$
the same $A$ is a rotation-scaling with no eigenvectors at
all, and only the passage to $\C$ reveals its normal form.
\end{example}

\begin{theorem}[Hermitian spectral theorem]\label{thm:b2:hermitian:spectral}
Every Hermitian endomorphism of a Hermitian space has an orthonormal
basis of eigenvectors (with real eigenvalues): $A^\dagger = A$
implies $A = U D U^{\dagger}$ with $U \in U(n)$ and $D$ real
diagonal.\index{spectral theorem!Hermitian}
\end{theorem}

\begin{proof}
Over $\C$, the characteristic polynomial splits: an eigenvector
$e_1$ exists (\cref{ch:b2:reduction}) --- no compactness argument
needed, one advantage of $\C$. Normalize it. Its orthogonal
complement $F = e_1^\perp$ is stable: for $x \perp e_1$,
\[
\langle e_1, u(x)\rangle = \langle u(e_1), x\rangle
= \lambda_1\langle e_1, x\rangle = 0
\]
($\lambda_1$ real). The restriction is Hermitian; induct on the
dimension and concatenate. In detail: the restriction $u|_F$ is
an endomorphism of the Hermitian space $F$ (dimension $n - 1$)
with $\langle u|_F(x), y\rangle = \langle x, u|_F(y)\rangle$
inherited from $u$; the induction hypothesis provides an
orthonormal basis $(e_2, \dots, e_n)$ of $F$ of eigenvectors,
and $(e_1, e_2, \dots, e_n)$ is orthonormal in $E$ ($e_1 \perp
F$) and made of eigenvectors of $u$. Matrix translation: the
columns of $U$ are the $e_i$, $U^\dagger U = I$ expresses their
orthonormality, and $AU = UD$ collects the eigenvalue
equations, whence $A = UDU^\dagger$ with $D$ real diagonal
(\cref{prop:b2:hermitian:eigenvalues}).
\end{proof}

\begin{example}\label{ex:b2:hermitian:pauli}
$A = \begin{pmatrix} 0 & -\iu\\ \iu & 0\end{pmatrix}$ is Hermitian
($A^\dagger = A$): eigenvalues from $\chi_A = X^2 - 1$: $\pm 1$
(real, as promised), with orthonormal eigenvectors
$\frac{1}{\sqrt2}(1, \iu)^{\mathsf T}$ and $\frac{1}{\sqrt2}(1,
-\iu)^{\mathsf T}$. (Physicists know $A$ as a Pauli matrix; the
reality of Hermitian spectra is why quantum observables are modeled
by Hermitian operators.)
\end{example}

\begin{example}[A positive definite Hermitian matrix, worked]\label{ex:b2:hermitian:pdrun}
$A = \begin{pmatrix} 2 & 1-\iu\\ 1+\iu & 3\end{pmatrix}$:
Hermitian, since the diagonal is real and the off-diagonal
entries are conjugates. Characteristic polynomial:
\[
(2-\lambda)(3-\lambda) - \abs{1-\iu}^2
= \lambda^2 - 5\lambda + 4
= (\lambda - 1)(\lambda - 4) :
\]
spectrum $\{1, 4\}$, real and positive --- $A$ is positive
definite. Eigenvectors: for $\lambda = 4$, the system $(A -
4I)v = 0$ gives $v_4 = (1 - \iu,\ 2)$ (check the second row:
$(1+\iu)(1-\iu) - 2 = 0$); for $\lambda = 1$, $v_1 = (1 - \iu,\
-1)$. Orthogonality, with the conjugate in the first slot:
\[
\langle v_4, v_1\rangle
= \conj{(1-\iu)}\,(1-\iu) + \conj{2}\,(-1)
= 2 - 2 = 0 . \checkmark
\]
Normalizing ($\norm{v_4}^2 = 2 + 4 = 6$, $\norm{v_1}^2 = 2 + 1
= 3$) gives the unitary $U = \bigl(\frac{v_4}{\sqrt6}\
\frac{v_1}{\sqrt3}\bigr)$ with $A =
U\operatorname{diag}(4,1)U^\dagger$. Closing insight: the
Rayleigh reading is immediate --- on the unit sphere of
$\C^2$, $\langle x, Ax\rangle$ ranges over $\intcc{1}{4}$,
attained at the two eigenvectors; this is the $n = 2$ germ of
the Courant--Fischer theory built in the weekend problem.
Spot-check that the form is real off the eigenvectors too: at
$x = (1, \iu)$,
\[
Ax = \bigl(2 + (1-\iu)\iu,\ (1+\iu) + 3\iu\bigr)
= (3 + \iu,\ 1 + 4\iu),
\]
\[
\langle x, Ax\rangle = \conj{1}\,(3+\iu) +
\conj{\iu}\,(1+4\iu)
= (3 + \iu) + (-\iu)(1 + 4\iu) = 3 + \iu - \iu + 4 = 7 \in
\R ,
\]
as \cref{prop:b2:hermitian:eigenvalues}'s proof mechanism
(conjugate-symmetry against $A^\dagger = A$) guarantees for
every $x$.
\end{example}

\begin{example}[A unitary matrix diagonalized]\label{ex:b2:hermitian:unitary}
$U = \frac{1}{\sqrt2}\begin{pmatrix} 1 & \iu\\ \iu & 1
\end{pmatrix}$ (unitary by \cref{exo:b2:hermitian:2}). Its
characteristic polynomial is $\bigl(X - \frac{1}{\sqrt2}\bigr)^2
+ \frac12$, with roots
\[
\lambda_\pm = \frac{1 \pm \iu}{\sqrt2} = \eu^{\pm\iu\pi/4},
\]
unimodular as \cref{prop:b2:hermitian:eigenvalues} promised, and
orthonormal eigenvectors $\frac{1}{\sqrt2}(1, \pm1)$. So $U =
V\operatorname{diag}(\eu^{\iu\pi/4},
\eu^{-\iu\pi/4})V^\dagger$: in the right basis, $U$ is a pair of
plane rotations by $\pm\frac\pi4$ --- a real rotation matrix has
no real eigenvectors, but over $\C$ it splits into two
unimodular scalars. Closing insight: Hermitian spectra live on
the real line, unitary spectra on the unit circle; both are
shadows of the same normality, and the Cayley transform of
\cref{exo:b2:hermitian:6} maps one picture to the other.
\end{example}

\begin{example}[The Cayley transform, computed]\label{ex:b2:hermitian:cayleyrun}
Run \cref{exo:b2:hermitian:6} on $H = \begin{pmatrix} 0 & 1\\
1 & 0\end{pmatrix}$ (Hermitian, spectrum $\{1, -1\}$,
orthonormal eigenvectors $\frac{1}{\sqrt2}(1, \pm1)$). In the
eigenbasis everything is scalar: the transform $\lambda
\mapsto \frac{\lambda - \iu}{\lambda + \iu}$ sends
\[
1 \longmapsto \frac{1 - \iu}{1 + \iu} = -\iu,
\qquad
-1 \longmapsto \frac{-1 - \iu}{-1 + \iu} = \iu ,
\]
(multiply by the conjugate of the denominator), so $U = (H -
\iu I)(H + \iu I)^{-1}$ is the unitary with eigenvalues
$\mp\iu$ on those same eigenvectors:
\[
U = \frac{1}{2}\begin{pmatrix} 1 & 1\\ 1 & -1\end{pmatrix}
\begin{pmatrix} -\iu & 0\\ 0 & \iu\end{pmatrix}
\begin{pmatrix} 1 & 1\\ 1 & -1\end{pmatrix}
= \begin{pmatrix} 0 & -\iu\\ -\iu & 0\end{pmatrix} .
\]
Check: $U^\dagger U = I$, and $1 \notin \operatorname{Sp}U =
\{\pm\iu\}$, as the theory promises. Closing insight: the real
line maps onto the unit circle minus the point $1$ ---
eigenvalue by eigenvalue, the Cayley transform is the
M\"obius map $\frac{\lambda-\iu}{\lambda+\iu}$, and matrices
just follow their spectra.
\end{example}

\begin{remark}[Common pitfalls]
\emph{(i) Where the bar falls:} this book conjugates the
\emph{first} slot, so coordinates are $\langle e_i, x\rangle$
and $\langle\lambda x, y\rangle = \conj\lambda\langle x,
y\rangle$; many texts conjugate the second slot instead ---
translate before comparing formulas, or signs of $\iu$ go
wrong silently. \emph{(ii) Complex polarization is stronger:}
over $\C$, if $\langle x, u(x)\rangle = 0$ for \emph{all} $x$
then $u = 0$ (expand $x + y$ and $x + \iu y$: both the real
and imaginary parts of $\langle x, u(y)\rangle$ vanish); over
$\R$ this fails --- the rotation by $\frac\pi2$ satisfies
$\langle x, u(x)\rangle = 0$ everywhere. Consequently, over
$\C$ only, ``$\langle x, u(x)\rangle \in \R$ for all $x$''
already forces $u$ Hermitian. \emph{(iii) Real normal is not
diagonalizable:} the matrix of
\cref{ex:b2:hermitian:skewexample} is normal but has no real
eigenvalue; unitary diagonalization is a theorem \emph{over}
$\C$, and over $\R$ one only gets block reductions.
\emph{(iv) Checking unitarity:} $U^\dagger U = I$ means the
\emph{columns} are orthonormal for the Hermitian product ---
testing $UU^{\mathsf T}$, or forgetting the conjugation in the
column products, are the two classic ways to certify a wrong
matrix.
\end{remark}

\begin{example}[Isometries are exactly the unitaries]\label{ex:b2:hermitian:isometry}
Norm preservation looks weaker than unitarity, but over $\C$ it
is not: if $\norm{u(x)} = \norm x$ for all $x$, then $u^*u =
\mathrm{id}$. Indeed $v = u^*u - \mathrm{id}$ is Hermitian and
satisfies $\langle x, v(x)\rangle = \norm{u(x)}^2 - \norm x^2 =
0$ for every $x$; by complex polarization (pitfall (ii)
above), a map with identically vanishing ``diagonal'' is zero:
$v = 0$. Concretely, the polarization runs
\[
0 = \langle x + y, v(x+y)\rangle
= \langle x, v(y)\rangle + \langle y, v(x)\rangle,
\qquad
0 = \langle x + \iu y, v(x + \iu y)\rangle
= \iu\langle x, v(y)\rangle - \iu\langle y, v(x)\rangle ,
\]
and the two lines together force $\langle x, v(y)\rangle = 0$
for all $x, y$. Closing insight: this is why ``unitary'' can
be checked by measuring lengths alone --- rigidity that the
Fourier chapter will exploit, where preserving the energy
$\norm f_2$ (Parseval) is the same as preserving all inner
products of coefficients.
\end{example}

\begin{remark}[Perspectives within this volume]
The Hermitian machinery built here is consumed almost
immediately. The Fourier chapter \emph{is} Hermitian geometry
in infinite dimension: the exponentials $(e_n)$ are an
orthonormal family for $\langle f, g\rangle =
\frac{1}{2\pi}\int\conj fg$, Bessel's inequality is the
projection estimate of this chapter's
\cref{thm:b2:hermitian:cs}, and Parseval is its limiting
equality. The finite Fourier transform
(\cref{exo:b2:hermitian:10}) reappears whenever convolution
must be diagonalized. And this chapter's weekend problem ---
Courant--Fischer, Weyl, interlacing --- supplies the
eigenvalue stability that the differential-equations chapter
invokes when it asserts that small perturbations of a system
move its frequencies only slightly. Backward, everything here
is the complex mirror of the quadratic-forms chapter: keep the
two dictionaries side by side ($A^{\mathsf T} \leftrightarrow
A^\dagger$, orthogonal $\leftrightarrow$ unitary, Rayleigh
real in both).
\end{remark}

\begin{remark}[Normal endomorphisms]
Over $\C$ the definitive statement is: $u$ is \emph{unitarily
diagonalizable if and only if it is normal} ($u^*u = uu^*$) ---
covering Hermitian, unitary and skew-Hermitian maps at once. The
proof is a pleasant strengthening of the argument above
(\cref{exo:b2:hermitian:8}). Over $\R$, by contrast, normality only
buys block-diagonalization (rotation blocks): complex geometry is
genuinely simpler.
\end{remark}

\begin{remark}[Where this is used]
Hermitian spectral theory is the mathematics of quantum
mechanics: observables are modeled by Hermitian operators
(real spectra = measurable values), time evolution by unitary
ones (norm preservation = conservation of probability). Within
this book, the Fourier chapter rests on the orthonormality of
the exponentials --- a Hermitian inner-product statement --- and
the diagonalization of circulant matrices
(\cref{exo:b2:hermitian:10}) is the finite Fourier transform.
The weekend problem develops the variational calculus of
eigenvalues (Courant--Fischer, Weyl, interlacing), the daily
bread of numerical analysis and mathematical physics; the Year
3 volume extends it to compact self-adjoint operators on
Hilbert spaces.
\end{remark}

\section{Exercises}

\begin{exercise}[$\star$]\label{exo:b2:hermitian:1}
On $\C^2$: compute $\langle x, y\rangle$, $\norm x$, $\norm y$ for
$x = (1, \iu)$, $y = (\iu, 1)$; are they orthogonal? Give an
orthonormal basis containing $\frac{x}{\norm x}$.
\end{exercise}

\begin{exercise}[$\star$]\label{exo:b2:hermitian:2}
Which are Hermitian? unitary? normal?
\[
\begin{pmatrix} 1 & \iu\\ -\iu & 2 \end{pmatrix},
\qquad
\frac{1}{\sqrt2}\begin{pmatrix} 1 & \iu\\ \iu & 1\end{pmatrix},
\qquad
\begin{pmatrix} 0 & 1\\ 0 & 0 \end{pmatrix}.
\]
\end{exercise}

\begin{exercise}[$\star$]\label{exo:b2:hermitian:3}
Prove that a matrix $A \in \mathcal{M}_n(\C)$ writes uniquely $A =
H + \iu K$ with $H, K$ Hermitian \emph{(the ``real and imaginary
parts'' $H = \frac{A + A^\dagger}{2}$, $K = \frac{A -
A^\dagger}{2\iu}$)}, and that $A$ is normal iff $H$ and $K$
commute.
\end{exercise}

\begin{exercise}[$\star\star$]\label{exo:b2:hermitian:4}
Diagonalize in an orthonormal basis: $A = \begin{pmatrix} 2 & \iu\\
-\iu & 2\end{pmatrix}$, and compute $A^k$ for $k \in \N$.
\end{exercise}

\begin{exercise}[$\star\star$]\label{exo:b2:hermitian:5}
Prove that $U(n)$ is compact, and that the eigenvalue map is
onto: every unimodular $\lambda$ arises for some unitary matrix.
Prove that $\det U \in \mathbb{U}$ (the unit circle) for $U \in
U(n)$.
\end{exercise}

\begin{exercise}[$\star\star$]\label{exo:b2:hermitian:6}
(Cayley transform) Let $H$ be Hermitian. Prove that $H + \iu I$ is
invertible and that $U = (H - \iu I)(H + \iu I)^{-1}$ is unitary,
with $1 \notin \operatorname{Sp}(U)$. \emph{(Work spectrally: on an
eigenbasis of $H$, everything is scalar.)}
\end{exercise}

\begin{exercise}[$\star\star$]\label{exo:b2:hermitian:7}
For $A$ Hermitian positive definite ($\langle x, Ax\rangle > 0$ for
$x \neq 0$), prove that $\operatorname{Sp}(A) \subseteq
\intoo{0}{\infty}$, that $A = B^2$ for a Hermitian positive
definite $B$, and that $\det A > 0$.
\end{exercise}

\begin{exercise}[$\star\star\star$]\label{exo:b2:hermitian:8}
(Spectral theorem for normal endomorphisms) Let $u$ be normal on a
Hermitian space.
\begin{enumerate}
  \item Prove $\norm{u(x)} = \norm{u^*(x)}$ for all $x$, and deduce
        $\ker(u - \lambda) = \ker(u^* - \conj\lambda)$.
  \item Prove that eigenspaces of $u$ for distinct eigenvalues are
        orthogonal, and that the orthogonal complement of an
        eigenspace is $u$-stable.
  \item Conclude by induction that $u$ is unitarily diagonalizable;
        and conversely.
\end{enumerate}
\end{exercise}

\begin{exercise}[$\star$]\label{exo:b2:hermitian:9}
An endomorphism is \emph{skew-Hermitian} when $u^* = -u$. Prove
that its eigenvalues are purely imaginary, that $u \mapsto \iu
u$ is a bijection from Hermitian to skew-Hermitian
endomorphisms, and that skew-Hermitian endomorphisms are
unitarily diagonalizable (\cref{exo:b2:hermitian:8}).
\end{exercise}

\begin{exercise}[$\star\star$]\label{exo:b2:hermitian:10}
(The finite Fourier transform) Let $S$ be the cyclic shift of
$\C^n$: $S(x_0, x_1, \dots, x_{n-1}) = (x_{n-1}, x_0, \dots,
x_{n-2})$, and $\omega = \eu^{2\iu\pi/n}$.
\begin{enumerate}
  \item Show that $S$ is unitary, and that the vectors $f_k =
        \frac{1}{\sqrt n}\bigl(1, \omega^k, \omega^{2k}, \dots,
        \omega^{(n-1)k}\bigr)$, $0 \leq k < n$, form an
        orthonormal basis of eigenvectors: $Sf_k = \omega^{-k}
        f_k$.
  \item Deduce that every \emph{circulant} matrix $C =
        \sum_{j=0}^{n-1} c_jS^j$ is normal, diagonalized by the
        same basis, with eigenvalues $\widehat c(k) = \sum_j
        c_j\,\omega^{-jk}$.
\end{enumerate}
\end{exercise}

\begin{exercise}[$\star\star$]\label{exo:b2:hermitian:11}
Let $P$ be an idempotent ($P^2 = P$) endomorphism of a Hermitian
space. Prove that $P$ is the \emph{orthogonal} projection onto
$\operatorname{im} P$ if and only if $P^* = P$. Give the matrix
of the orthogonal projection onto $\C v$ ($\norm v = 1$), and
onto a subspace with orthonormal basis $(v_1, \dots, v_k)$.
\end{exercise}

\begin{exercise}[$\star\star\star$]\label{exo:b2:hermitian:12}
(Spectral projectors by interpolation) Let $A$ be Hermitian with
\emph{distinct} eigenvalues $\lambda_1, \dots, \lambda_p$ and
eigenspace decomposition $E = \bigoplus_i E_i$. Define the
Lagrange polynomials $L_i(X) = \prod_{j\neq i}\frac{X -
\lambda_j}{\lambda_i - \lambda_j}$. Prove that $P_i = L_i(A)$
is the orthogonal projection onto $E_i$, that $P_iP_j = 0$ for
$i \neq j$, $\sum_i P_i = I$, and $A = \sum_i \lambda_iP_i$
(the \emph{spectral decomposition}); express $f(A)$ for any
polynomial $f$ in terms of the $P_i$.
\end{exercise}

\section{Problem: Courant--Fischer, Weyl, and the calculus of
eigenvalues}

\begin{problem}\label{pb:b2:hermitian:1}
The eigenvalues of a Hermitian matrix are not just roots of a
polynomial: they are solutions of \emph{optimization problems}.
That variational point of view --- Rayleigh quotients and the
\emph{Courant--Fischer min-max theorem}\index{Courant--Fischer
theorem} --- makes eigenvalues comparable, stable, and
computable, and this problem harvests its classical crops:
\emph{Weyl's perturbation inequalities}\index{Weyl's
inequality}, \emph{Cauchy interlacing}, the Schur and Ky Fan
trace inequalities, the monotonicity of the matrix square root,
and the spectrum of the discrete Laplacian. Throughout, $A, B,
E$ are Hermitian on $E = \C^n$ with eigenvalues listed in
decreasing order $\lambda_1(A) \geq \dots \geq \lambda_n(A)$,
and $R_A(x) = \frac{\langle x, Ax\rangle}{\langle x, x\rangle}$
for $x \neq 0$ is the \emph{Rayleigh quotient}\index{Rayleigh
quotient}.

\medskip
\noindent\textbf{Part I --- Rayleigh quotients and min-max.}
Fix an orthonormal eigenbasis $(e_1, \dots, e_n)$, $Ae_i =
\lambda_ie_i$.
\begin{enumerate}
  \item Show that $R_A(x)$ is real, and that
        \[
        \lambda_n \leq R_A(x) \leq \lambda_1
        \qquad (x \neq 0),
        \]
        both bounds attained: $\lambda_1 = \max R_A$,
        $\lambda_n = \min R_A$.
  \item Show that the critical points of $R_A$ are exactly the
        eigenvectors of $A$ \emph{(expand $t \mapsto R_A(x +
        tv)$ at $t = 0$ for $v$ arbitrary, then replace $v$ by
        $\iu v$)}.
  \item Let $V_k = \operatorname{Vect}(e_1, \dots, e_k)$ and
        $W_k = \operatorname{Vect}(e_k, \dots, e_n)$. Show
        \[
        \min_{x \in V_k\setminus\{0\}} R_A(x) = \lambda_k
        = \max_{x \in W_k\setminus\{0\}} R_A(x) .
        \]
  \item Prove the \emph{Courant--Fischer theorem}: for $1 \leq
        k \leq n$,
        \[
        \lambda_k = \max_{\dim V = k}\;\min_{x \in
        V\setminus\{0\}} R_A(x)
        = \min_{\dim W = n-k+1}\;\max_{x \in
        W\setminus\{0\}} R_A(x)
        \]
        \emph{(for any $V$ of dimension $k$: $V \cap W_k \neq
        \{0\}$ by Grassmann, so $\min_V R_A \leq \lambda_k$;
        question 3 shows the bound is attained)}.
  \item (Monotonicity) Write $A \leq B$ when $B - A$ is
        positive semidefinite. Deduce from question 4: $A \leq
        B$ implies $\lambda_k(A) \leq \lambda_k(B)$ for every
        $k$.
\end{enumerate}

\medskip
\noindent\textbf{Part II --- Weyl's inequalities.}
\begin{enumerate}[resume]
  \item Show that subspaces $V, W \subseteq \C^n$ with $\dim V
        + \dim W > n$ intersect nontrivially, and generalize:
        $\dim(V_1 \cap V_2 \cap V_3) \geq \dim V_1 + \dim V_2 +
        \dim V_3 - 2n$.
  \item Prove \emph{Weyl's inequality}: for $i + j - 1 \leq n$,
        \[
        \lambda_{i+j-1}(A + B) \leq \lambda_i(A) +
        \lambda_j(B)
        \]
        \emph{(intersect the subspaces $W_i(A)$, $W_j(B)$ and
        $V_{i+j-1}(A+B)$ of question 3 and count
        dimensions)}.
  \item Define $\vertiii{E}_2 = \max_{\norm x = 1}\norm{Ex}$
        and show $\vertiii E_2 = \max_k\abs{\lambda_k(E)}$ for
        Hermitian $E$. Deduce \emph{Weyl's perturbation
        theorem}:
        \[
        \bigl|\lambda_k(A + E) - \lambda_k(A)\bigr| \leq
        \vertiii{E}_2
        \qquad (1 \leq k \leq n) :
        \]
        each eigenvalue is a $1$-Lipschitz function of the
        matrix.
  \item (Rank-one perturbations) Let $P$ be Hermitian positive
        semidefinite of rank $1$. Show
        \[
        \lambda_k(A) \leq \lambda_k(A + P) \leq
        \lambda_{k-1}(A) \qquad (2 \leq k \leq n),
        \]
        together with $\lambda_1(A) \leq \lambda_1(A+P)$: the
        new eigenvalues \emph{interlace} the old ones.
  \item Check question 8 numerically: $A = \begin{pmatrix} 2 &
        \iu\\ -\iu & 2\end{pmatrix}$
        (\cref{exo:b2:hermitian:4}: spectrum $\{3, 1\}$) and
        $E = \begin{pmatrix} 0 & 1\\ 1 & 0\end{pmatrix}$
        (spectrum $\{1, -1\}$): compute the spectrum of $A +
        E$ and both sides of the inequality.
\end{enumerate}

\medskip
\noindent\textbf{Part III --- Interlacing and trace
inequalities.}
\begin{enumerate}[resume]
  \item (Cauchy interlacing) Let $B$ be the leading $(n-1)
        \times(n-1)$ principal submatrix of $A$. Prove
        \[
        \lambda_{k+1}(A) \leq \lambda_k(B) \leq \lambda_k(A)
        \qquad (1 \leq k \leq n-1)
        \]
        \emph{(view $\C^{n-1} \subseteq \C^n$; on it, $R_B$ is
        the restriction of $R_A$; apply Courant--Fischer on
        both levels)}.
  \item Iterate: for a principal submatrix $B$ of size $n - m$,
        $\lambda_{k+m}(A) \leq \lambda_k(B) \leq
        \lambda_k(A)$.
  \item (Schur) Let $d_1 \geq d_2 \geq \dots \geq d_n$ be the
        diagonal entries of $A$, sorted. Prove, for every $k$:
        \[
        \sum_{i=1}^{k} d_i \leq \sum_{i=1}^{k}\lambda_i(A),
        \]
        with equality at $k = n$ (the trace)
        \emph{(the $k$ chosen diagonal entries form a
        principal $k\times k$ submatrix; bound its trace by
        question 12)}.
  \item (Ky Fan) Prove
        \[
        \sum_{i=1}^{k}\lambda_i(A)
        = \max\Bigl\{\sum_{i=1}^{k}\langle x_i, Ax_i\rangle
        : (x_1, \dots, x_k) \text{ orthonormal}\Bigr\} .
        \]
  \item Verify questions 11 and 13 on
        \[
        A = \begin{pmatrix} 2 & 1 & 0\\ 1 & 2 & 1\\
        0 & 1 & 2\end{pmatrix}
        \qquad
        \bigl(\operatorname{Sp} = \{2 - \sqrt2,\ 2,\
        2+\sqrt2\}\bigr)
        \]
        against its leading $2\times2$ block (spectrum $\{1,
        3\}$) and its diagonal.
\end{enumerate}

\medskip
\noindent\textbf{Part IV --- The Loewner order.} $A \leq B$
still means $B - A$ positive semidefinite; all matrices in this
part are Hermitian.
\begin{enumerate}[resume]
  \item Show: $A \leq B$ implies $a_{ii} \leq b_{ii}$ for all
        $i$, $\operatorname{tr} A \leq \operatorname{tr} B$,
        and $C^\dagger AC \leq C^\dagger BC$ for \emph{every}
        complex matrix $C$.
  \item Show that squaring is \emph{not} monotone: for
        \[
        A = \begin{pmatrix} 1 & 0\\ 0 & 0\end{pmatrix},
        \qquad
        B = \begin{pmatrix} 2 & 1\\ 1 & 1\end{pmatrix},
        \]
        check $0 \leq A \leq B$ but $A^2 \not\leq B^2$.
  \item Prove that the square root \emph{is} monotone: $0 \leq
        A \leq B$ implies $\sqrt A \leq \sqrt B$ \emph{(let
        $\mu$ be an eigenvalue of $\sqrt B - \sqrt A$ with unit
        eigenvector $v$; compute $\langle v, (B - A)v\rangle =
        \mu\bigl(\langle v, \sqrt B\,v\rangle + \langle v,
        \sqrt A\,v\rangle\bigr)$ and discuss)}.
  \item Prove that inversion is antitone on positive definite
        matrices: $0 < A \leq B$ implies $B^{-1} \leq A^{-1}$
        \emph{(congruate by $A^{-1/2}$ to reduce to $I \leq M
        \Rightarrow M^{-1} \leq I$, which is scalar in a
        spectral basis)}.
  \item Let $A, B$ be positive definite. Show that the
        eigenvalues of $AB$ (not Hermitian in general!) are
        real and positive, and that
        \[
        \lambda_{\max}(AB) \leq
        \lambda_{\max}(A)\,\lambda_{\max}(B)
        \]
        \emph{(conjugate by $\sqrt A$: $AB \sim \sqrt
        A\,B\sqrt A$)}.
\end{enumerate}

\medskip
\noindent\textbf{Part V --- The discrete Laplacian, worked.}
Let $T_n$ be the $n \times n$ tridiagonal matrix with $2$ on
the diagonal and $-1$ on the two adjacent diagonals.
\begin{enumerate}[resume]
  \item With $\theta_k = \frac{k\pi}{n+1}$, verify that the
        vectors $v_k = \bigl(\sin(j\theta_k)\bigr)_{1\leq j\leq
        n}$ satisfy $T_nv_k = (2 - 2\cos\theta_k)\,v_k$
        \emph{(product-to-sum identity; check the boundary rows
        $j = 1, n$)}. Conclude:
        \[
        \operatorname{Sp}(T_n) = \Bigl\{4\sin^2
        \frac{k\pi}{2(n+1)} : 1 \leq k \leq n\Bigr\},
        \]
        all simple, all positive: $T_n$ is positive definite.
  \item (A potential) For a real diagonal $D =
        \operatorname{diag}(d_1, \dots, d_n)$, sandwich the
        spectrum: for every $k$,
        \[
        \lambda_k(T_n) + \min_i d_i \;\leq\;
        \lambda_k(T_n + D) \;\leq\; \lambda_k(T_n) + \max_i
        d_i .
        \]
  \item Check Cauchy interlacing between $T_3$ and $T_2$
        explicitly (spectra $\{2 \pm \sqrt2, 2\}$ and $\{1,
        3\}$), and interpret: $T_2$ is $T_3$ with one endpoint
        of the path removed.
  \item Show the extreme eigenvalues satisfy, as $n \to
        \infty$:
        \[
        \lambda_{\min}(T_n) = 4\sin^2\frac{\pi}{2(n+1)}
        \sim \frac{\pi^2}{(n+1)^2},
        \qquad
        \lambda_{\max}(T_n) \to 4 ,
        \]
        so the condition number $\kappa_n =
        \lambda_{\max}/\lambda_{\min}$ grows like
        $\frac{4(n+1)^2}{\pi^2}$: discretizing a second
        derivative on a finer and finer grid is
        intrinsically ill-conditioned.
  \item Synthesis. In one sentence each: (i) why the
        variational characterization, not the characteristic
        polynomial, is what makes eigenvalues stable
        (questions 8--9); (ii) which questions used only
        $\lambda_1 = \max R_A$ and which needed the full
        min-max; (iii) what the Loewner order adds to the
        story; (iv) where these tools reappear (numerical
        analysis of question 24's stiffness matrices; quantum
        perturbation theory; and, in the Year 3 volume, the
        min-max principle for compact self-adjoint
        operators).
\end{enumerate}
\end{problem}
